\begin{textsample}{POS Dim 1 – sp – Score 49.00 – sp/sp\_em\_43.txt}  \label{ex:f1_pos_111}
Itinerário Formativo - Área de Ciências Humanas e Sociais Aplicadas Os itinerários formativos têm como objetivo consolidar e aprofundar conhecimentos, preparar o estudante para os desafios do mundo do trabalho e da cidadania na \textbf{contemporaneidade} e \textbf{aprimorar} a formação ética, além de promover uma postura ativa frente ao conhecimento científico, filosófico e à produção artística e literária.
Ao estruturar os itinerários formativos para a área de Ciências Humanas e Sociais Aplicadas, buscou -se nas \textbf{competências} de área descritas na BNCC referência para as habilidades elencadas para o itinerário de Ciências Humanas, conforme Portaria nº 1.432, de dezembro de 2018, que estabelece os referenciais para elaboração dos itinerários formativos.
A Resolução nº 3, de 21 de novembro de 2018, que atualiza as DCNEM, foi outro aporte fundamental para a proposição dos itinerários.
O formato apresentado reconhece e valoriza a liberdade, a autonomia e a responsabilidade das unidades escolares quanto à concepção, formulação e execução de suas propostas pedagógicas, assim como o \textbf{potencial} de cada unidade escolar em organizar e atualizar os espaços escolares para dinamizar e diversificar o ensino e a aprendizagem.
A estrutura dos itinerários formativos orienta para a valorização dos adolescentes, jovens e adultos, assim como seus tempos de aprendizagem.
Logo, optou -se por considerar nos “ pressupostos \textbf{metodológicos} ” orientações gerais, em que objetivos fundamentais se integram às esferas da pesquisa, do trabalho e das práticas sociais e \textbf{socioambientais}, tendo como referência as habilidades propostas pelos eixos estruturantes.
Considerando a responsabilidade que se estabelece pela educação formal de preparar o estudante para o mundo \textbf{contemporâneo}, que \textbf{exige} cidadãos \textbf{críticos}, \textbf{reflexivos}, éticos, flexíveis e abertos para o novo ou para a renovação do que é conhecido, e mediante o dinamismo da sociedade \textbf{contemporânea}, é \textbf{imperativo} não separar forma e conteúdo.
Ou seja, as metodologias devem estar de acordo com os objetivos e estes com os objetos de conhecimento.
Destaca -se, nesse sentido, a centralidade da “ situação-problema ”, que deve ser pensada e reconhecida não apenas como uma intervenção para resolver questões de forma \textbf{imediata}, mas também a \textbf{médio} e longo prazo.
A situação-problema não pode ser reduzida a uma questão de reparo ; ela deve ser considerada também no âmbito da atualização e \textbf{contextualização} das aprendizagens, isto é, precisa possibilitar mais do que a observação, a interpretação e a intervenção, para melhor se apropriar de questões do universo humano, em busca de um mundo melhor.
Portanto, \textbf{espera} -se que os docentes considerem as metodologias ativas como \textbf{caminhos} para oportunizar o desenvolvimento das habilidades previstas no itinerário formativo da área de Ciências Humanas e Sociais Aplicadas e nos itinerários integrados desta área com as demais.
Nesse sentido, a resolução de \textbf{problemas} reais e atividades que \textbf{demandam} pesquisas, leituras, organização para produção e compartilhamento de conhecimento, tendo em \textbf{vista} as \textbf{demandas} do mundo \textbf{contemporâneo}. \textbf{Esperam} -se dos professores mais flexibilidade e abertura para o novo em relação aos ritmos de aprendizagem, ao trabalho com materiais impressos e \textbf{digitais}, à organização e orientação de atividades.
Eles também devem ser capazes de incentivar o estudante a ser produtor e não apenas receptor do conhecimento, assim como rever as formas de acompanhamento das aprendizagens.
INVESTIGAÇÃO CIENTÍFICA A investigação científica deverá apresentar oportunidades de realização dentro e fora da sala de \textbf{aula}, em diferentes espaços.
O estudo de campo e a visita dirigida a museus, parques, centros culturais, bibliotecas especializadas, entre outros locais, devem ser considerados tendo em \textbf{vista} os interesses do estudante em relação ao tema pesquisado.
O registro merece atenção especial no que se refere ao seu desenvolvimento, à ampliação do vocabulário, ao incremento na \textbf{descrição} dos dados, entre outros elementos reveladores do desenvolvimento do estudante na pesquisa científica no campo das Ciências Humanas e Sociais Aplicadas.
Dessa forma, neste eixo, além das referências bibliográficas, será fundamental \textbf{explicitar} para o estudante como são construídas e debatidas as questões de pesquisa nas Ciências Humanas e Sociais Aplicadas.
Logo, contribuições da tradição dos componentes curriculares da área devem ser exploradas nesse contexto.
Considera -se também que, no desenvolvimento deste eixo, o estudante possa conhecer e experimentar diferentes percursos da investigação a partir do referencial : \textbf{abordagem} quantitativa e/ou \textbf{qualitativa} ; natureza ( pesquisa básica e pesquisa aplicada ) e objetivos ( pesquisa exploratória, descritiva, experimental, bibliográfica, documental, de meio, estudo de caso, pesquisa participante, pesquisa-ação, pesquisa etnográfica, entre outras \textbf{pertinentes}, considerando a natureza e a \textbf{abordagem} do \textbf{problema} ).
PROCESSOS CRIATIVOS Os processos criativos devem provocar a prática, isto é, não se trata apenas de propor ideias, mas também de reconhecer e ressignificar soluções e percursos para chegar a conclusões e/ou propor projetos socioeconômicos, \textbf{socioambientais}, de preservação da memória e da cultura, entre outros.
Nesse sentido, os componentes da área de Ciências Humanas e Sociais Aplicadas devem estar abertos para a utilização e produção de imagens, filmes, redes sociais etc., para encaminhar processos na escola e/ou na comunidade.
No contexto deste eixo, deve -se abrir espaço para a livre manifestação de ideias de todos os envolvidos, de forma a estimular o \textbf{debate}, a tomada de posição, a defesa de um ponto de \textbf{vista} por meio de \textbf{argumentos} e a possibilidade de mudar de posição diante de \textbf{argumentos} mais consistentes.
A pesquisa deve ser o ponto \textbf{central} para fundamentar uma ideia, um procedimento, assim como a percepção quanto à realização de uma \textbf{tarefa} e o desenvolvimento de um projeto.
A organização, no contexto deste eixo, permite não apenas a divisão de \textbf{tarefas}, mas também o compartilhamento de responsabilidades acerca das atividades a serem desenvolvidas, os recursos a serem utilizados e os registros que comportam a avaliação, autoavaliação e avaliação entre pares, em todas as etapas do processo a ser desenvolvido sob a inspiração deste eixo.
No âmbito do compartilhamento de responsabilidades, entende -se que o eixo possibilita, ainda, reflexões acerca da tomada de decisão, inclusive, sobre mudanças de rumo, quando se mostrarem necessárias.
Todos esses movimentos devem ser acompanhados e avaliados como parte da educação integral do estudante.
MEDIAÇÃO E INTERVENÇÃO SOCIOCULTURAL Este eixo relaciona -se com a reflexão apoiada em dados oficiais, escuta da comunidade, entre outras ações \textbf{pertinentes}.
No contexto da área de Ciências Humanas e Sociais Aplicadas, este eixo está diretamente relacionado com o eixo Investigação Científica.
De onde \textbf{decorre} a necessidade de acompanhamento, uma vez que devem se acordar protocolos para a coleta de dados e para a escuta \textbf{qualificada}, tendo em \textbf{vista} a ação de mediação e intervenção sociocultural.
A observação desses protocolos \textbf{configura} -se em um desafio para professores, equipe gestora e estudantes, especialmente, considerando -se a possibilidade de estar diante de diferentes territorialidades e questões socioculturais e \textbf{socioambientais}.
O processo de mediação e intervenção requer observação, registro e o constante \textbf{debate} acerca dos valores presentes em propostas de colaboração e mediação, assim como a identificação do \textbf{problema}, suas características territoriais e temporais, a formulação de hipóteses, a fundamentação \textbf{teórica}, a decisão sobre o cronograma de estudo/aprendizagem a necessidade de ajustes, avaliação e implementação da ação.
O eixo Mediação e Intervenção Sociocultural deverá aprofundar habilidades relacionadas à observação, interpretação e proposição de resolução de situações-problema.
EMPREENDEDORISMO O eixo Empreendedorismo envolve -se com a \textbf{mobilização} de conhecimentos e o reconhecimento do que ainda precisa ser desenvolvido para empreender projetos pessoais e produtivos que estejam no \textbf{horizonte} do projeto de vida do estudante.
Especificamente no campo das Ciências Humanas e Sociais Aplicadas, destaca -se a \textbf{análise} dos \textbf{fundamentos} da ética, de forma a conhecer e propor empreendimentos capazes de valorizar a liberdade, a cooperação, a autonomia e a convivência democrática, entre outros valores.
É preciso considerar, no contexto deste eixo, que o estudante deve assumir uma postura participante, atuante e engajada, exercitando a disposição de pensar e repensar, tendo em \textbf{vista} a sustentabilidade dos seus projetos e o \textbf{potencial} \textbf{transformador} de suas ações.
Assim, aprendizagens sobre valores humanos, simbólicos e criativos, o valor patrimonial da cultura brasileira e sua dimensão histórica, social, artística e ambiental, entre outros, podem ser inspiradoras para a efetivação deste eixo.
Por fim, é preciso levar o estudante a compreender que o projeto de vida não se resume a uma escolha pessoal, mas envolve o exercício da cidadania, da empatia, da \textbf{capacidade} de entender o mundo e reconhecer -se como parte dele.
Isso \textbf{posto}, ao \textbf{final} da Educação Básica, \textbf{espera} -se que o estudante do itinerário da área de Ciências Humanas e Sociais Aplicadas seja capaz de : 1.
Apreciar as diferentes manifestações da experiência do pensamento, com aportes ao exercício da suspeita às respostas fáceis, de forma a elaborar perguntas e \textbf{argumentos} visando à proposição de soluções éticas, estéticas e criativas para as diferentes questões do mundo \textbf{contemporâneo}. 2.
Compreender e interpretar conceitos de natureza, paisagem, território, rede, cultura, fenômeno social e globalização nas dinâmicas históricas, \textbf{econômicas}, sociais e políticas, reconhecendo seus \textbf{desdobramentos} nas representações e produções culturais. 3.
Atribuir sentido e significado para a vida em sociedade por meio de permanências e regularidades, e também mudanças e transformações, no processo histórico, tendo como referência a dimensão da temporalidade e do trabalho como expressão da organização da vida material, individual e coletiva e das relações de poder \textbf{explicitadas} nos processos de dominação, hegemonia, resistência e participação política. 4.
Perceber a complexidade das relações que cada um estabelece consigo e com a comunidade em que vive, interferindo na sua identidade e na diversidade, material e simbólica, de forma a ser capaz de mediar relações com criticidade e ética, atuando para a superação das condições individualista, egocêntrica, preconceituosa, racista e hedonista, entre outras, em seus diferentes vieses na sociedade \textbf{contemporânea}.
Seja nas relações com a natureza e/ou entre pessoas, utilizar habilidades \textbf{conceituais} e técnicas para \textbf{mobilizar} diferentes \textbf{atores} no empreendimento de ações socioculturais e \textbf{socioambientais} em âmbito local, regional, nacional e/ou global. 5.
Vivenciar conscientemente o processo cognitivo \textbf{conceitual}, de forma a reconhecer -se como produtor que \textbf{mobiliza} conhecimentos e práticas na constituição de um projeto de vida viável e sustentável. 6.
Reconhecer que o projeto de vida é estratégico e envolve escolhas pessoais, relacionando -as com as \textbf{demandas} da sociedade \textbf{contemporânea}.
Isso requer compreender que a história de cada um, seus interesses e seus sonhos devem ser considerados também na relação com o contexto social, político e cultural em que se vive, \textbf{exigindo} \textbf{análise} \textbf{crítica} das limitações sociais e das desigualdades, bem como flexibilidade e motivação para a aprendizagem ao longo da vida.
A seguir encontra -se o \textbf{organizador} curricular do itinerário formativo da área de Ciências Humanas e Sociais Aplicadas que contempla as habilidades relacionadas às \textbf{competências} gerais da Educação Básica, as habilidades específicas dos itinerários formativos dessa área do conhecimento e os pressupostos \textbf{metodológicos} que visam orientar e indicar possibilidades para a concretização das aprendizagens \textbf{esperadas}, para cada eixo estruturante.

% matched lemmas: abordagem, análise, aprimorar, argumento, ator, aula, caminho, capacidade, central, competência, conceitual, configurar, contemporaneidade, contemporâneo, contextualização, crítico, debate, decorrer, demanda, demandar, descrição, desdobramento, digital, econômico, esperar, exigir, explicitar, final, fundamento, horizonte, imediato, imperativo, metodológico, mobilizar, mobilização, médio, organizador, pertinente, postar, potencial, problema, qualificar, qualitativo, reflexivo, socioambiental, tarefa, teórico, transformador, vista
\end{textsample}
